% Source PDF pages 107--109; manual pages 2-78--2-80.
\section{Guidance Rules}\label{sec:guidance-rules}

Because TAOS is a three-degree-of-freedom point-mass trajectory simulation, the
body attitude or angular orientation cannot be obtained from the equations of
motion. Instead, body attitude is supplied through guidance rules. TAOS assumes
that the vehicle has a control system capable of maintaining the commanded
attitude.

The body-fixed coordinate system is defined relative to another coordinate system
by three angles. Relative to the geocentric, geodetic, or inertial-platform systems,
these are three Euler angles. Relative to a velocity coordinate system, they are a
pair of aerodynamic angles and a bank angle. The valid combinations are listed in
\cref{tab:body-attitude-control-variables}.

% Several earlier semantic longtables are intentionally unnumbered but still advance
% LaTeX's table counter. Restore the source-manual sequence before Tables 2-2 and 2-3.
\setcounter{table}{1}

\begin{table}[htbp]
  \centering
  \caption{Body-attitude control variables.}
  \label{tab:body-attitude-control-variables}
  \small
  \renewcommand{\arraystretch}{1.12}
  \begin{tabularx}{\linewidth}{@{}c>{\raggedright\arraybackslash}X>{\raggedright\arraybackslash}X>{\raggedright\arraybackslash}X@{}}
    \toprule
    Set & \multicolumn{3}{c}{Body-attitude angles}\\
    \midrule
    1 & $\alpha$ Angle of attack & $\beta_E$ Euler sideslip angle & $\mu_{gc}$ Geocentric bank angle\\
    2 & $\alpha$ Angle of attack & $\beta$ Sideslip angle & $\mu_{gc}$ Geocentric bank angle\\
    3 & $\alpha_T$ Total angle of attack & $\phi_w$ Windward meridian & $\mu_{gc}$ Geocentric bank angle\\
    4 & $\alpha$ Angle of attack & $\beta_E$ Euler sideslip angle & $\mu_{gd}$ Geodetic bank angle\\
    5 & $\alpha$ Angle of attack & $\beta$ Sideslip angle & $\mu_{gd}$ Geodetic bank angle\\
    6 & $\alpha_T$ Total angle of attack & $\phi_w$ Windward meridian & $\mu_{gd}$ Geodetic bank angle\\
    7 & $\Psi_{gc}$ Geocentric yaw angle & $\Theta_{gc}$ Geocentric pitch angle & $\Phi_{gc}$ Geocentric roll angle\\
    8 & $\Psi_{gd}$ Geodetic yaw angle & $\Theta_{gd}$ Geodetic pitch angle & $\Phi_{gd}$ Geodetic roll angle\\
    9 & $\Psi_p$ Inertial yaw angle & $\Theta_p$ Inertial pitch angle & $\Phi_p$ Inertial roll angle\\
    \bottomrule
  \end{tabularx}
\end{table}

\taossubhead{Control Variables}

Vehicle acceleration, and therefore vehicle motion, is determined by the forces
acting on the vehicle. These forces generally depend on aerodynamic angles, which
in turn depend on the body-attitude angles. The attitude angles therefore control the
forces that control the trajectory and can be regarded as \emph{control variables}.

For a vehicle of fixed geometry, aerodynamic forces commonly depend on altitude,
velocity, angle of attack, and sideslip angle. Propulsive forces commonly depend on
altitude, velocity, and power setting. Altitude and velocity follow from the vehicle
state, but angle of attack, sideslip, and power setting are independent quantities that
must be supplied to integrate the equations of motion.

TAOS uses four control variables. Three are selected from one row of
\cref{tab:body-attitude-control-variables}; the fourth is power setting $P$. A
control variable affects the trajectory only when the corresponding force tables
depend on it. For example, if the propulsion tables are not functions of power
setting, changing $P$ has no effect on propulsive force or trajectory.

\taossubhead{Guidance Rules}

Guidance rules determine the control-variable values. A rule may specify a control
variable directly or may specify a flight condition from which TAOS solves for a
control variable. Four control variables are used, so four guidance rules are required,
although controls whose desired value is zero may use their default.

A direct rule such as ``fly at a $5^\circ$ angle of attack'' supplies a control value
without further calculation. Directly supplied attitude variables must form one of
the consistent sets in \cref{tab:body-attitude-control-variables}; all directly
specified angles must come from the same row.

An indirect rule such as ``fly at a $0^\circ$ vertical flight-path angle'' commands
level flight. The required angle of attack is not known in advance, so TAOS solves
for the control value that produces the requested condition. The indirect rules and
their associated solution variables are listed in
\cref{tab:indirect-guidance-rules}.

\begin{table}[htbp]
  \centering
  \caption{Guidance rules that indirectly specify control variables.}
  \label{tab:indirect-guidance-rules}
  \small
  \renewcommand{\arraystretch}{1.08}
  \begin{tabularx}{\linewidth}{@{}>{\raggedright\arraybackslash}X>{\raggedright\arraybackslash}p{0.25\textwidth}>{\centering\arraybackslash}p{0.20\textwidth}@{}}
    \toprule
    Guidance rule & Control variables & Solution variable\\
    \midrule
    Fly at a constant altitude & $\alpha,\alpha_T,P$ & $\ddot h$\\
    Fly at a constant lift coefficient & $\alpha,\alpha_T$ & $C_L$\\
    Fly at a constant side-force coefficient & $\beta,\beta_E,\phi_w$ & $C_S$\\
    Fly at a constant dynamic pressure & $P,\alpha,\alpha_T$ & $\ddot q$\\
    Fly at a constant geocentric flight-path angle & $\alpha,\alpha_T,P$ & $\dot\gamma_{gc}$\\
    Fly at a constant geodetic flight-path angle & $\alpha,\alpha_T,P$ & $\dot\gamma_{gd}$\\
    Fly intercept or proportional navigation in pitch & $\alpha,\alpha_T$ & $\dot\gamma_{gd}$\\
    Fly intercept or proportional navigation in yaw & $\beta,\beta_E,\phi_w,\mu_{gc},\mu_{gd}$ & $\dot\psi_{gd}$\\
    Fly at a constant lift-to-drag ratio & $\alpha,\alpha_T$ & $L/D$\\
    Fly at the maximum lift-to-drag ratio & $\alpha,\alpha_T$ & $\alpha_{L/D}$\\
    \bottomrule
  \end{tabularx}
\end{table}

\begin{table}[htbp]
  \ContinuedFloat
  \centering
  \caption[]{Guidance rules that indirectly specify control variables (continued).}
  \small
  \renewcommand{\arraystretch}{1.08}
  \begin{tabularx}{\linewidth}{@{}>{\raggedright\arraybackslash}X>{\raggedright\arraybackslash}p{0.25\textwidth}>{\centering\arraybackslash}p{0.20\textwidth}@{}}
    \toprule
    Guidance rule & Control variables & Solution variable\\
    \midrule
    Fly at a constant Mach number & $P,\alpha,\alpha_T$ & $\dot M$\\
    Fly at a constant axial-$g$ loading & $P,\alpha,\alpha_T$ & $\vec a_{sp}\cdot\hat x_b$\\
    Fly at a constant lateral-$g$ loading & $\beta,\beta_E,\phi_w$ & $\vec a_{sp}\cdot\hat y_b$\\
    Fly at a constant normal-$g$ loading & $\alpha,\alpha_T$ & $\vec a_{sp}\cdot\hat z_b$\\
    Fly at a constant geocentric heading angle & $\beta,\beta_E,\phi_w,\mu_{gc},\mu_{gd}$ & $\dot\psi_{gc}$\\
    Fly at a constant geodetic heading angle & $\beta,\beta_E,\phi_w,\mu_{gc},\mu_{gd}$ & $\dot\psi_{gd}$\\
    Fly at a constant thrust & $P$ & $\lVert\vec F_{\mathrm{prop}}\rVert$\\
    Fly up/down the range-insensitive axis in pitch & $\alpha,\alpha_T$ & $r_{\mathrm{iip}}$\\
    Fly up/down the range-insensitive axis in yaw & $\beta,\beta_E,\phi_w,\mu_{gc},\mu_{gd}$ & $\alpha_{\mathrm{iip}}$\\
    Fly at a constant velocity & $P,\alpha,\alpha_T$ & $\lVert\vec V_{\oplus}\rVert$\\
    \bottomrule
  \end{tabularx}
\end{table}

A solution variable is a function of the control variables and is associated with its
guidance rule. It is generally an acceleration-like quantity that can change
instantaneously along the trajectory.
